% Part: incompleteness
% Chapter: representability-in-q
% Section: composition-representable

\documentclass[../../../include/open-logic-section]{subfiles}

\begin{document}

\olfileid{inc}{req}{cmp}
\olsection{التركيب قابل للتمثيل في $\Th{Q}$}

لتكن $h$ هي الدالة التي يعرّفها التركيب
\[
h(x_0,\dots,x_{l-1}) = f(g_0(x_0,\dots,x_{l-1}), \dots,
g_{k-1}(x_0,\dots,x_{l-1})).
\]
وقد عُيّنت من قبل !!{formula}s $!A_f,!A_{g_0},\dots,!A_{g_{k-1}}$
لتمثيل الدوال $f$ و$g_0$، \dots، و$g_{k-1}$، بالترتيب.
فالمطلوب الآن !!a{formula}~$!A_h$ تمثل~$h$.

نأخذ أولًا أيسر الحالات، وهي أن تكون الدوال كلها أحادية الموضع:
$h(x)=f(g(x))$. ولتكن $!A_f(y,z)$ ممثلةً لـ$f$،
و$!A_g(x,y)$ ممثلةً لـ$g$؛ فنطلب !!a{formula}~$!A_h(x,z)$
لتمثيل~$h$. ومفتاح البناء أن $h(x)=z$ تكافئ وجود~$y$
تصدق معها $z=f(y)$ و$y=g(x)$ معًا. فمتى كانت $h(x)=z$،
كانت $g(x)$ هي الشاهد~$y$؛ ومتى وجد مثل هذا $y$،
أنتجت $y=g(x)$ مع $z=f(y)$ المساواة $z=f(g(x))$.
ولهذا نختار $\lexists[y][(!A_g(x,y)\land !A_f(y,z))]$
لتكون $!A_h(x,z)$، ثم نتحقق من أن $\Th{Q}$ تثبت
ما يلزم من !!{formula}s.

\begin{prop}
\ollabel{prop:rep1}
متى كان $h(n)=m$ لزم $\Th{Q}\Proves !A_h(\num{n},\num{m})$.
\end{prop}

\begin{proof}
ليكن $h(n)=m$، أي $f(g(n))=m$، ولنسمّ القيمة الوسطى $k=g(n)$.
فنحصل على ما يأتي:
\begin{align*}
  \Th{Q} & \Proves !A_g(\num{n}, \num{k})
  \intertext{لأن $!A_g$ تمثل~$g$، و}
  \Th{Q} & \Proves !A_f(\num{k}, \num{m})
  \intertext{لأن $!A_f$ تمثل~$f$. ومن ثم}
  \Th{Q} & \Proves !A_g(\num{n}, \num{k}) \land !A_f(\num{k}, \num{m})
  \intertext{وبالتالي أيضًا}
  \Th{Q} & \Proves \lexists[y][(!A_g(\num{n}, y) \land !A_f(y, \num{m}))],
\end{align*}
وهذه هي بالضبط $\Th{Q}\Proves !A_h(\num{n},\num{m})$.
\end{proof}

\begin{prop}
\ollabel{prop:rep2}
متى كان $h(n)=m$ لزم $\Th{Q}\Proves\lforall[z][(!A_h(\num{n},z)\lif
z=\num{m})]$.
\end{prop}

\begin{proof}
نفرض $h(n)=m$، وهو $f(g(n))=m$، ونضع $k=g(n)$.
ومن شرط انفراد القيمة في التمثيل نستخرج:
\begin{align*}
  \Th{Q} & \Proves \lforall[y][(!A_g(\num{n}, y) \lif \eq[y][\num{k}])]
  \intertext{لأن $!A_g$ تمثل~$g$، و}
  \Th{Q} & \Proves \lforall[z][(!A_f(\num{k}, z) \lif \eq[z][\num{m}])]
  \intertext{لأن $!A_f$ تمثل~$f$. وبقليل من المنطق يمكننا أن نبيّن أيضًا أن}
  \Th{Q} & \Proves \lforall[z][(\lexists[y][(!A_g(\num{n}, y) \land
      !A_f(y, z))] \lif \eq[z][\num{m}])].
\end{align*}
وهذا، مع تغيير اسم المتغير المربوط، هو $\Th{Q}\Proves\lforall[y][(!A_h(\num n,y)\lif
\eq[y][\num m])]$.
\end{proof}

وعلى هذه الطريقة يجري البرهان في الحالة العامة، حين تكون رتبتا
$f$ و$g_i$ أكبر من~$1$، وإن طالت الكتابة.

\begin{prop}
\ollabel{prop:rep-composition}
لتكن $!A_f(y_0,\dots,y_{k-1},z)$ ممثلة للدالة
$f(y_0,\dots,y_{k-1})$ في~$\Th{Q}$، ولتكن
$!A_{g_i}(x_0,\dots,x_{l-1},y)$ ممثلة للدالة
$g_i(x_0,\dots,x_{l-1})$ في~$\Th{Q}$ لكل $i<k$.
% تصحيح صياغة كلاسيكية (C295-01): عُيّن مجال الكم بالمؤشر i<k بدل لفظ «رتبة» الملتبس.
عندئذ تكون الصيغة
\begin{multline*}
  \lexists[y_0\dots][\lexists[y_{k-1}][(!A_{g_0}(x_0,\dots,x_{l-1},y_0) \land
      \dots \land {}]]\\
  !A_{g_{k-1}}(x_0,\dots,x_{l-1},y_{k-1}) \land !A_f(y_0,\dots,y_{k-1},z))
\end{multline*}
تمثيلًا للدالة
\[
h(x_0, \dots, x_{l-1}) = f(g_0(x_0, \dots, x_{l-1}), \dots, g_{k-1}(x_0,
\dots, x_{l-1})).
\]
\end{prop}

\begin{proof}
يُترك تفصيله تمرينًا.
\end{proof}

\begin{prob}
اتخذ برهانَي \olref[inc][req][cmp]{prop:rep1} و
\olref[inc][req][cmp]{prop:rep2} مثالًا، ثم فصّل برهان
\olref[inc][req][cmp]{prop:rep-composition}.
\end{prob}

\end{document}
